% --- Archon physics formalization source begin ---
% archon:physics
% archon:covers PhyXMiniProblems/problem_phyx_mini_0802.lean
% archon:source-report reports/phyx_mini/problem_phyx_mini_0802.source.json

\chapter{Physics problem phyx\_mini\_0802}
\label{ch:PhyXMiniProblems_problem_phyx_mini_0802}

\paragraph{Problem source.}
\#\# Physical scenario

An elevator and its load have a combined mass of \textbackslash{}( 800 \textbackslash{}mathrm\{kg\} \textbackslash{}) as shown in figure. The elevator is initially moving downward at \textbackslash{}( 10.0 \textbackslash{}mathrm\{m/s\} \textbackslash{}); it slows to a stop with constant acceleration in a distance of \textbackslash{}( 25.0 \textbackslash{}mathrm\{m\} \textbackslash{}).

\#\# Question

What is the tension \textbackslash{}( T \textbackslash{}) in the supporting cable while the elevator is being brought to rest?

\#\# Answer choices

- A: A: 8760N

- B: B: 9440N

- C: C: 9380N

- D: D: 8540N

\#\# Figure caption (auxiliary; use the image as primary evidence)

The image shows an elevator scenario on the left and a force diagram on the right. 

**Elevator Scene:**
- The elevator is enclosed within a shaft.
- It is depicted with an arrow pointing downward, labeled "Moving down with decreasing speed," indicating deceleration.
- The elevator doors are visible, along with an interface panel with buttons.

**Force Diagram:**
- The diagram illustrates a vertical axis with two main forces.
- The upward force is labeled "T" (tension).
- The downward force is labeled "w = mg" (weight).
- There is an upward arrow labeled "a\_y" indicating acceleration, pointing opposite to the direction of movement, consistent with deceleration.
  
The elements convey that the elevator is slowing down as it moves downward due to the net force being upward.

\#\# Dataset metadata

Dynamics; Physical Model Grounding Reasoning, Multi-Formula Reasoning

\paragraph{Recorded answer/context.}
B: B: 9440N

\paragraph{Figure/image path.}
/root/proposal\_for\_physic/science-mango/phyx\_mini\_run/phyx\_data/test\_image/802.png

\paragraph{Formalization target.}
create a compiling Lean file with sorry bodies at `PhyXMiniProblems/problem\_phyx\_mini\_0802.lean`. The Lean declarations must preserve the physical quantities, dimensions or dimensional roles, figure labels, governing-law hypotheses, and final relation expressed by this problem.
Use Mathlib/Physlib names found through LeanExplore where available. If a physics API is missing, introduce faithful local abstractions rather than scalar placeholder aliases.

\begin{theorem}[Physics formalization target]
\label{thm:physics:phyx_mini_0802:target}
\lean{PhyXMiniProblems.ProblemPhyXMini0802.elevatorCableTension_is_answerB}
\uses{def:physics:phyx-mini-0802:phyxminiproblems-problemphyxmini0802-forceinnewtons, def:physics:phyx-mini-0802:phyxminiproblems-problemphyxmini0802-elevatorbrakingsetup, def:physics:phyx-mini-0802:phyxminiproblems-problemphyxmini0802-matchesproblemdescription, def:physics:phyx-mini-0802:phyxminiproblems-problemphyxmini0802-matchesprimaryfigure, def:physics:phyx-mini-0802:phyxminiproblems-problemphyxmini0802-satisfiesconstantaccelerationkinematics, def:physics:phyx-mini-0802:phyxminiproblems-problemphyxmini0802-satisfiesverticaldynamics, def:physics:phyx-mini-0802:phyxminiproblems-problemphyxmini0802-recordedanswerchoice, def:physics:phyx-mini-0802:phyxminiproblems-problemphyxmini0802-matchesanswerchoice}
The assigned autoformalize agent should translate this physics problem into Lean declarations in the covered file, with theorem and lemma proof bodies written as `by sorry`.
\end{theorem}
\begin{proof}
This is an autoformalization task, not a proof task. Produce statements that can later be proved by the physics prover without weakening the physical model.
\end{proof}

% --- Archon declaration topology begin ---
\section{Declaration topology}
\begin{definition}[velocity Dimension]
  \label{def:physics:phyx-mini-0802:phyxminiproblems-problemphyxmini0802-velocitydimension}
  \lean{PhyXMiniProblems.ProblemPhyXMini0802.velocityDimension}
  The physical dimension of velocity, L T\textasciicircum{}-1
\end{definition}
\begin{proof}
  This is the declared model definition of velocity Dimension.
\end{proof}
\begin{definition}[acceleration Dimension]
  \label{def:physics:phyx-mini-0802:phyxminiproblems-problemphyxmini0802-accelerationdimension}
  \lean{PhyXMiniProblems.ProblemPhyXMini0802.accelerationDimension}
  The physical dimension of acceleration, L T\textasciicircum{}-2
\end{definition}
\begin{proof}
  This is the declared model definition of acceleration Dimension.
\end{proof}
\begin{definition}[force Dimension]
  \label{def:physics:phyx-mini-0802:phyxminiproblems-problemphyxmini0802-forcedimension}
  \lean{PhyXMiniProblems.ProblemPhyXMini0802.forceDimension}
  The physical dimension of force, M L T\textasciicircum{}-2
\end{definition}
\begin{proof}
  This is the declared model definition of force Dimension.
\end{proof}
\begin{definition}[Mass Quantity]
  \label{def:physics:phyx-mini-0802:phyxminiproblems-problemphyxmini0802-massquantity}
  \lean{PhyXMiniProblems.ProblemPhyXMini0802.MassQuantity}
  A nonnegative physical mass
\end{definition}
\begin{proof}
  This is the declared model definition of Mass Quantity.
\end{proof}
\begin{definition}[Length Magnitude Quantity]
  \label{def:physics:phyx-mini-0802:phyxminiproblems-problemphyxmini0802-lengthmagnitudequantity}
  \lean{PhyXMiniProblems.ProblemPhyXMini0802.LengthMagnitudeQuantity}
  A nonnegative physical length, used for the distance travelled
\end{definition}
\begin{proof}
  This is the declared model definition of Length Magnitude Quantity.
\end{proof}
\begin{definition}[Vertical Displacement Quantity]
  \label{def:physics:phyx-mini-0802:phyxminiproblems-problemphyxmini0802-verticaldisplacementquantity}
  \lean{PhyXMiniProblems.ProblemPhyXMini0802.VerticalDisplacementQuantity}
  A signed component of displacement along the pictured Y axis
\end{definition}
\begin{proof}
  This is the declared model definition of Vertical Displacement Quantity.
\end{proof}
\begin{definition}[Speed Magnitude Quantity]
  \label{def:physics:phyx-mini-0802:phyxminiproblems-problemphyxmini0802-speedmagnitudequantity}
  \lean{PhyXMiniProblems.ProblemPhyXMini0802.SpeedMagnitudeQuantity}
  \uses{def:physics:phyx-mini-0802:phyxminiproblems-problemphyxmini0802-velocitydimension}
  A nonnegative physical speed
\end{definition}
\begin{proof}
  This is the declared model definition of Speed Magnitude Quantity.
\end{proof}
\begin{definition}[Vertical Velocity Quantity]
  \label{def:physics:phyx-mini-0802:phyxminiproblems-problemphyxmini0802-verticalvelocityquantity}
  \lean{PhyXMiniProblems.ProblemPhyXMini0802.VerticalVelocityQuantity}
  \uses{def:physics:phyx-mini-0802:phyxminiproblems-problemphyxmini0802-velocitydimension}
  A signed component of velocity along the pictured Y axis
\end{definition}
\begin{proof}
  This is the declared model definition of Vertical Velocity Quantity.
\end{proof}
\begin{definition}[Vertical Acceleration Quantity]
  \label{def:physics:phyx-mini-0802:phyxminiproblems-problemphyxmini0802-verticalaccelerationquantity}
  \lean{PhyXMiniProblems.ProblemPhyXMini0802.VerticalAccelerationQuantity}
  \uses{def:physics:phyx-mini-0802:phyxminiproblems-problemphyxmini0802-accelerationdimension}
  A signed component of acceleration along the pictured Y axis
\end{definition}
\begin{proof}
  This is the declared model definition of Vertical Acceleration Quantity.
\end{proof}
\begin{definition}[Acceleration Magnitude Quantity]
  \label{def:physics:phyx-mini-0802:phyxminiproblems-problemphyxmini0802-accelerationmagnitudequantity}
  \lean{PhyXMiniProblems.ProblemPhyXMini0802.AccelerationMagnitudeQuantity}
  \uses{def:physics:phyx-mini-0802:phyxminiproblems-problemphyxmini0802-accelerationdimension}
  A nonnegative magnitude of gravitational acceleration
\end{definition}
\begin{proof}
  This is the declared model definition of Acceleration Magnitude Quantity.
\end{proof}
\begin{definition}[Force Magnitude Quantity]
  \label{def:physics:phyx-mini-0802:phyxminiproblems-problemphyxmini0802-forcemagnitudequantity}
  \lean{PhyXMiniProblems.ProblemPhyXMini0802.ForceMagnitudeQuantity}
  \uses{def:physics:phyx-mini-0802:phyxminiproblems-problemphyxmini0802-forcedimension}
  A nonnegative force magnitude, used for both tension and weight
\end{definition}
\begin{proof}
  This is the declared model definition of Force Magnitude Quantity.
\end{proof}
\begin{definition}[nonnegative Readout]
  \label{def:physics:phyx-mini-0802:phyxminiproblems-problemphyxmini0802-nonnegativereadout}
  \lean{PhyXMiniProblems.ProblemPhyXMini0802.nonnegativeReadout}
  Read a nonnegative dimensionful quantity in coherent units
\end{definition}
\begin{proof}
  This is the declared model definition of nonnegative Readout.
\end{proof}
\begin{definition}[signed Readout]
  \label{def:physics:phyx-mini-0802:phyxminiproblems-problemphyxmini0802-signedreadout}
  \lean{PhyXMiniProblems.ProblemPhyXMini0802.signedReadout}
  Read a signed dimensionful quantity in coherent units
\end{definition}
\begin{proof}
  This is the declared model definition of signed Readout.
\end{proof}
\begin{definition}[mass In Kilograms]
  \label{def:physics:phyx-mini-0802:phyxminiproblems-problemphyxmini0802-massinkilograms}
  \lean{PhyXMiniProblems.ProblemPhyXMini0802.massInKilograms}
  \uses{def:physics:phyx-mini-0802:phyxminiproblems-problemphyxmini0802-massquantity, def:physics:phyx-mini-0802:phyxminiproblems-problemphyxmini0802-nonnegativereadout}
  Kilogram readout of a physical mass
\end{definition}
\begin{proof}
  This is the declared model definition of mass In Kilograms.
\end{proof}
\begin{definition}[distance In Meters]
  \label{def:physics:phyx-mini-0802:phyxminiproblems-problemphyxmini0802-distanceinmeters}
  \lean{PhyXMiniProblems.ProblemPhyXMini0802.distanceInMeters}
  \uses{def:physics:phyx-mini-0802:phyxminiproblems-problemphyxmini0802-lengthmagnitudequantity, def:physics:phyx-mini-0802:phyxminiproblems-problemphyxmini0802-nonnegativereadout}
  Metre readout of a nonnegative distance
\end{definition}
\begin{proof}
  This is the declared model definition of distance In Meters.
\end{proof}
\begin{definition}[displacement YIn Meters]
  \label{def:physics:phyx-mini-0802:phyxminiproblems-problemphyxmini0802-displacementyinmeters}
  \lean{PhyXMiniProblems.ProblemPhyXMini0802.displacementYInMeters}
  \uses{def:physics:phyx-mini-0802:phyxminiproblems-problemphyxmini0802-verticaldisplacementquantity, def:physics:phyx-mini-0802:phyxminiproblems-problemphyxmini0802-signedreadout}
  Metre readout of a signed vertical displacement
\end{definition}
\begin{proof}
  This is the declared model definition of displacement YIn Meters.
\end{proof}
\begin{definition}[speed In Meters Per Second]
  \label{def:physics:phyx-mini-0802:phyxminiproblems-problemphyxmini0802-speedinmeterspersecond}
  \lean{PhyXMiniProblems.ProblemPhyXMini0802.speedInMetersPerSecond}
  \uses{def:physics:phyx-mini-0802:phyxminiproblems-problemphyxmini0802-speedmagnitudequantity, def:physics:phyx-mini-0802:phyxminiproblems-problemphyxmini0802-nonnegativereadout}
  Metres-per-second readout of a speed magnitude
\end{definition}
\begin{proof}
  This is the declared model definition of speed In Meters Per Second.
\end{proof}
\begin{definition}[velocity YIn Meters Per Second]
  \label{def:physics:phyx-mini-0802:phyxminiproblems-problemphyxmini0802-velocityyinmeterspersecond}
  \lean{PhyXMiniProblems.ProblemPhyXMini0802.velocityYInMetersPerSecond}
  \uses{def:physics:phyx-mini-0802:phyxminiproblems-problemphyxmini0802-verticalvelocityquantity, def:physics:phyx-mini-0802:phyxminiproblems-problemphyxmini0802-signedreadout}
  Metres-per-second readout of a signed vertical velocity component
\end{definition}
\begin{proof}
  This is the declared model definition of velocity YIn Meters Per Second.
\end{proof}
\begin{definition}[acceleration YIn Meters Per Second Squared]
  \label{def:physics:phyx-mini-0802:phyxminiproblems-problemphyxmini0802-accelerationyinmeterspersecondsquared}
  \lean{PhyXMiniProblems.ProblemPhyXMini0802.accelerationYInMetersPerSecondSquared}
  \uses{def:physics:phyx-mini-0802:phyxminiproblems-problemphyxmini0802-verticalaccelerationquantity, def:physics:phyx-mini-0802:phyxminiproblems-problemphyxmini0802-signedreadout}
  Metres-per-second-squared readout of signed vertical acceleration
\end{definition}
\begin{proof}
  This is the declared model definition of acceleration YIn Meters Per Second Squared.
\end{proof}
\begin{definition}[acceleration Magnitude In Meters Per Second Squared]
  \label{def:physics:phyx-mini-0802:phyxminiproblems-problemphyxmini0802-accelerationmagnitudeinmeterspersecondsquared}
  \lean{PhyXMiniProblems.ProblemPhyXMini0802.accelerationMagnitudeInMetersPerSecondSquared}
  \uses{def:physics:phyx-mini-0802:phyxminiproblems-problemphyxmini0802-accelerationmagnitudequantity, def:physics:phyx-mini-0802:phyxminiproblems-problemphyxmini0802-nonnegativereadout}
  Metres-per-second-squared readout of an acceleration magnitude
\end{definition}
\begin{proof}
  This is the declared model definition of acceleration Magnitude In Meters Per Second Squared.
\end{proof}
\begin{definition}[force In Newtons]
  \label{def:physics:phyx-mini-0802:phyxminiproblems-problemphyxmini0802-forceinnewtons}
  \lean{PhyXMiniProblems.ProblemPhyXMini0802.forceInNewtons}
  \uses{def:physics:phyx-mini-0802:phyxminiproblems-problemphyxmini0802-forcemagnitudequantity, def:physics:phyx-mini-0802:phyxminiproblems-problemphyxmini0802-nonnegativereadout}
  Newton readout of a force magnitude
\end{definition}
\begin{proof}
  This is the declared model definition of force In Newtons.
\end{proof}
\begin{definition}[Vertical Direction]
  \label{def:physics:phyx-mini-0802:phyxminiproblems-problemphyxmini0802-verticaldirection}
  \lean{PhyXMiniProblems.ProblemPhyXMini0802.VerticalDirection}
  Qualitative vertical directions used by the primary figure
\end{definition}
\begin{proof}
  This is the declared model definition of Vertical Direction.
\end{proof}
\begin{definition}[Speed Trend]
  \label{def:physics:phyx-mini-0802:phyxminiproblems-problemphyxmini0802-speedtrend}
  \lean{PhyXMiniProblems.ProblemPhyXMini0802.SpeedTrend}
  Whether the magnitude of the elevator's velocity is increasing or decreasing
\end{definition}
\begin{proof}
  This is the declared model definition of Speed Trend.
\end{proof}
\begin{definition}[Mechanical System Kind]
  \label{def:physics:phyx-mini-0802:phyxminiproblems-problemphyxmini0802-mechanicalsystemkind}
  \lean{PhyXMiniProblems.ProblemPhyXMini0802.MechanicalSystemKind}
  The object whose combined mass is specified in the problem
\end{definition}
\begin{proof}
  This is the declared model definition of Mechanical System Kind.
\end{proof}
\begin{definition}[Acceleration Regime]
  \label{def:physics:phyx-mini-0802:phyxminiproblems-problemphyxmini0802-accelerationregime}
  \lean{PhyXMiniProblems.ProblemPhyXMini0802.AccelerationRegime}
  The acceleration regime stated in the problem
\end{definition}
\begin{proof}
  This is the declared model definition of Acceleration Regime.
\end{proof}
\begin{definition}[Elevator Figure]
  \label{def:physics:phyx-mini-0802:phyxminiproblems-problemphyxmini0802-elevatorfigure}
  \lean{PhyXMiniProblems.ProblemPhyXMini0802.ElevatorFigure}
  \uses{def:physics:phyx-mini-0802:phyxminiproblems-problemphyxmini0802-verticaldirection, def:physics:phyx-mini-0802:phyxminiproblems-problemphyxmini0802-speedtrend}
  Literal categorical information retained from the primary raster. The question's unknown tension magnitude does not occur as a numerical figure field
\end{definition}
\begin{proof}
  This is the declared model definition of Elevator Figure.
\end{proof}
\begin{definition}[Elevator Braking Setup]
  \label{def:physics:phyx-mini-0802:phyxminiproblems-problemphyxmini0802-elevatorbrakingsetup}
  \lean{PhyXMiniProblems.ProblemPhyXMini0802.ElevatorBrakingSetup}
  \uses{def:physics:phyx-mini-0802:phyxminiproblems-problemphyxmini0802-massquantity, def:physics:phyx-mini-0802:phyxminiproblems-problemphyxmini0802-lengthmagnitudequantity, def:physics:phyx-mini-0802:phyxminiproblems-problemphyxmini0802-verticaldisplacementquantity, def:physics:phyx-mini-0802:phyxminiproblems-problemphyxmini0802-speedmagnitudequantity, def:physics:phyx-mini-0802:phyxminiproblems-problemphyxmini0802-verticalvelocityquantity, def:physics:phyx-mini-0802:phyxminiproblems-problemphyxmini0802-verticalaccelerationquantity, def:physics:phyx-mini-0802:phyxminiproblems-problemphyxmini0802-accelerationmagnitudequantity, def:physics:phyx-mini-0802:phyxminiproblems-problemphyxmini0802-forcemagnitudequantity, def:physics:phyx-mini-0802:phyxminiproblems-problemphyxmini0802-mechanicalsystemkind, def:physics:phyx-mini-0802:phyxminiproblems-problemphyxmini0802-accelerationregime, def:physics:phyx-mini-0802:phyxminiproblems-problemphyxmini0802-elevatorfigure}
  Independent physical quantities for the stopping interval. In particular, cableTension is not defined from an answer choice or from 9440
\end{definition}
\begin{proof}
  This is the declared model definition of Elevator Braking Setup.
\end{proof}
\begin{definition}[Matches Problem Description]
  \label{def:physics:phyx-mini-0802:phyxminiproblems-problemphyxmini0802-matchesproblemdescription}
  \lean{PhyXMiniProblems.ProblemPhyXMini0802.MatchesProblemDescription}
  \uses{def:physics:phyx-mini-0802:phyxminiproblems-problemphyxmini0802-massinkilograms, def:physics:phyx-mini-0802:phyxminiproblems-problemphyxmini0802-distanceinmeters, def:physics:phyx-mini-0802:phyxminiproblems-problemphyxmini0802-displacementyinmeters, def:physics:phyx-mini-0802:phyxminiproblems-problemphyxmini0802-speedinmeterspersecond, def:physics:phyx-mini-0802:phyxminiproblems-problemphyxmini0802-velocityyinmeterspersecond, def:physics:phyx-mini-0802:phyxminiproblems-problemphyxmini0802-accelerationmagnitudeinmeterspersecondsquared, def:physics:phyx-mini-0802:phyxminiproblems-problemphyxmini0802-elevatorbrakingsetup}
  The prose data expressed as SI readouts. Upward is positive, so the initial velocity and stopping displacement are the negatives of their stated magnitudes. The standard near-Earth value g = 9.8 m/s\textasciicircum{}2 is the convention used by the recorded answer choices
\end{definition}
\begin{proof}
  This is the declared model definition of Matches Problem Description.
\end{proof}
\begin{definition}[Matches Primary Figure]
  \label{def:physics:phyx-mini-0802:phyxminiproblems-problemphyxmini0802-matchesprimaryfigure}
  \lean{PhyXMiniProblems.ProblemPhyXMini0802.MatchesPrimaryFigure}
  \uses{def:physics:phyx-mini-0802:phyxminiproblems-problemphyxmini0802-elevatorbrakingsetup}
  Primary-image evidence. The label w = mg is recorded here only as a graphical fact; the corresponding physical law is stated independently in SatisfiesVerticalDynamics
\end{definition}
\begin{proof}
  This is the declared model definition of Matches Primary Figure.
\end{proof}
\begin{definition}[Satisfies Constant Acceleration Kinematics]
  \label{def:physics:phyx-mini-0802:phyxminiproblems-problemphyxmini0802-satisfiesconstantaccelerationkinematics}
  \lean{PhyXMiniProblems.ProblemPhyXMini0802.SatisfiesConstantAccelerationKinematics}
  \uses{def:physics:phyx-mini-0802:phyxminiproblems-problemphyxmini0802-signedreadout, def:physics:phyx-mini-0802:phyxminiproblems-problemphyxmini0802-elevatorbrakingsetup}
  The constant-acceleration relation v\_f\textasciicircum{}2 = v\_i\textasciicircum{}2 + 2 a\_y Δy, stated in every coherent unit system. It contains no numerical value for the unknown acceleration or tension
\end{definition}
\begin{proof}
  This is the declared model definition of Satisfies Constant Acceleration Kinematics.
\end{proof}
\begin{definition}[Satisfies Vertical Dynamics]
  \label{def:physics:phyx-mini-0802:phyxminiproblems-problemphyxmini0802-satisfiesverticaldynamics}
  \lean{PhyXMiniProblems.ProblemPhyXMini0802.SatisfiesVerticalDynamics}
  \uses{def:physics:phyx-mini-0802:phyxminiproblems-problemphyxmini0802-nonnegativereadout, def:physics:phyx-mini-0802:phyxminiproblems-problemphyxmini0802-signedreadout, def:physics:phyx-mini-0802:phyxminiproblems-problemphyxmini0802-elevatorbrakingsetup}
  The downward weight has magnitude mg, while the upward-positive force balance is T - w = m a\_y. These general dynamics laws contain neither the requested tension nor any answer-choice value
\end{definition}
\begin{proof}
  This is the declared model definition of Satisfies Vertical Dynamics.
\end{proof}
\begin{lemma}[braking Acceleration Y eq two]
  \label{lem:physics:phyx-mini-0802:phyxminiproblems-problemphyxmini0802-brakingaccelerationy-eq-two}
  \lean{PhyXMiniProblems.ProblemPhyXMini0802.brakingAccelerationY_eq_two}
  \uses{def:physics:phyx-mini-0802:phyxminiproblems-problemphyxmini0802-accelerationyinmeterspersecondsquared, def:physics:phyx-mini-0802:phyxminiproblems-problemphyxmini0802-elevatorbrakingsetup, def:physics:phyx-mini-0802:phyxminiproblems-problemphyxmini0802-matchesproblemdescription, def:physics:phyx-mini-0802:phyxminiproblems-problemphyxmini0802-satisfiesconstantaccelerationkinematics}
  The stopping data and constant-acceleration law imply a\_y = 2 m/s\textasciicircum{}2
\end{lemma}
\begin{proof}
  The conclusion follows from the governing relations and figure readouts named in the statement of braking Acceleration Y eq two.
\end{proof}
\begin{lemma}[elevator Weight eq 7840]
  \label{lem:physics:phyx-mini-0802:phyxminiproblems-problemphyxmini0802-elevatorweight-eq-7840}
  \lean{PhyXMiniProblems.ProblemPhyXMini0802.elevatorWeight_eq_7840}
  \uses{def:physics:phyx-mini-0802:phyxminiproblems-problemphyxmini0802-forceinnewtons, def:physics:phyx-mini-0802:phyxminiproblems-problemphyxmini0802-elevatorbrakingsetup, def:physics:phyx-mini-0802:phyxminiproblems-problemphyxmini0802-matchesproblemdescription, def:physics:phyx-mini-0802:phyxminiproblems-problemphyxmini0802-satisfiesverticaldynamics}
  The stated mass and standard gravity give a weight magnitude of 7840 N
\end{lemma}
\begin{proof}
  The conclusion follows from the governing relations and figure readouts named in the statement of elevator Weight eq 7840.
\end{proof}
\begin{definition}[Answer Choice]
  \label{def:physics:phyx-mini-0802:phyxminiproblems-problemphyxmini0802-answerchoice}
  \lean{PhyXMiniProblems.ProblemPhyXMini0802.AnswerChoice}
  Labels of the four answers printed in the problem
\end{definition}
\begin{proof}
  This is the declared model definition of Answer Choice.
\end{proof}
\begin{definition}[tension In Newtons]
  \label{def:physics:phyx-mini-0802:phyxminiproblems-problemphyxmini0802-answerchoice-tensioninnewtons}
  \lean{PhyXMiniProblems.ProblemPhyXMini0802.AnswerChoice.tensionInNewtons}
  \uses{def:physics:phyx-mini-0802:phyxminiproblems-problemphyxmini0802-answerchoice}
  Cable-tension value printed beside each answer label, in newtons
\end{definition}
\begin{proof}
  This is the declared model definition of tension In Newtons.
\end{proof}
\begin{definition}[recorded Answer Choice]
  \label{def:physics:phyx-mini-0802:phyxminiproblems-problemphyxmini0802-recordedanswerchoice}
  \lean{PhyXMiniProblems.ProblemPhyXMini0802.recordedAnswerChoice}
  \uses{def:physics:phyx-mini-0802:phyxminiproblems-problemphyxmini0802-answerchoice}
  The answer label recorded in the supplied dataset metadata
\end{definition}
\begin{proof}
  This is the declared model definition of recorded Answer Choice.
\end{proof}
\begin{definition}[Matches Answer Choice]
  \label{def:physics:phyx-mini-0802:phyxminiproblems-problemphyxmini0802-matchesanswerchoice}
  \lean{PhyXMiniProblems.ProblemPhyXMini0802.MatchesAnswerChoice}
  \uses{def:physics:phyx-mini-0802:phyxminiproblems-problemphyxmini0802-forceinnewtons, def:physics:phyx-mini-0802:phyxminiproblems-problemphyxmini0802-elevatorbrakingsetup, def:physics:phyx-mini-0802:phyxminiproblems-problemphyxmini0802-answerchoice}
  A displayed answer agrees exactly with the modeled cable-tension readout
\end{definition}
\begin{proof}
  This is the declared model definition of Matches Answer Choice.
\end{proof}
% --- Archon declaration topology end ---
% --- Archon physics formalization source end ---
